
\section{maincontent}

\begin{frame}
\frametitle{Strong-Field Ionization Phenomena Revealed by Quantum Trajectories \\
\small (Phys. Rev. Lett. \textbf{133}, 073201 (2024))}
\framesubtitle{Surrounding concepts}


\footnotesize
%\begin{textblock*}{200pt}(250pt,60pt)
%\begin{itemize}
%\item 3 steps model
%\includegraphics[scale=.20]{images/Color-online-Schematic-representation-of-the-dominant-physical-mechanisms-behind.png}
%\end{itemize}
%\end{textblock*}
 
\begin{textblock*}{400pt}(30pt,60pt)
\begin{itemize}
\item Carrier Envelope Phase
%%
\includegraphics[scale=.75]{images/polcls.png}
\begin{align}
\begin{split}
 \boldsymbol{F}(t) = F_0 f(t) \sin \left( \omega t + \varphi\right)
   \boldsymbol{\hat{e}_x}, \hspace{100pt} 
\boldsymbol{k} = \frac{\omega}{c}\boldsymbol{\hat{e}_z}
\end{split}
 \nonumber 
\end{align}
\item CEP effects
\begin{itemize}
\footnotesize
 \item[$\bullet$] Important parameters in the case of few cycles pulses
$(n_{oc} < 10 )$
 \item[$\bullet$] Directly effects most of the relevant observables
 \item[$\bullet$] Made evident by a privileged direction 
in the specra (anisotropy parameter)
 \item[$\bullet$] In the case of Ti:Sapphire, $I \sim 10^{14}\,{\rm W/cm^2}$,
most CEP effects gone by $T\geqslant 10\tau$
\end{itemize}
\item Energy level shifts
\begin{itemize}
\footnotesize
 \item[$\bullet$] Under the field, the fundamental drops 
by $-\alpha_0I/4$, with $\alpha_0$ the polarisability of the atom
 \item[$\bullet$] On the other hand, 
the continuum threshold increases with $U_p$
 \item[$\bullet$] For strong fields, Stark shift also contribute to the 
energy level shifts

\end{itemize}
%\item Excitation rates constructive interference condition
%\begin{align}
%\begin{split}
% \Delta\theta = T\left(  U_p  + I_p  +  E_{nl}  \right), \hspace{20pt}
%\Longrightarrow \hspace{20pt} m\omega = I_p  + E_{nl} + U_P  \hspace{100pt}
%\end{split}
% \nonumber 
%\end{align}

\end{itemize}
\end{textblock*}


\begin{textblock*}{100pt}(330pt,210pt)
\begin{align}
\footnotesize
U_p = \frac{F_0^2}{4\omega^2}, \hspace{20pt} 
\gamma = \sqrt{\frac{I_p}{2U_p}}\nonumber
\end{align}
\end{textblock*}

\end{frame}
%
%
%
\begin{frame}
\frametitle{Carrier-Envelope Phase-Dependent Strong-Field Excitation \\
\small (Phys. Rev. Lett. \textbf{128}, 173201 (2022))}
\framesubtitle{What to expect from this paper}

\small
\begin{textblock*}{400pt}(30pt,60pt)
\begin{itemize}
\item Few studies to date have explored the how CEP affects 
bound-state populations
\item CEP can be used to clearly determine when the dynamic changed from 
Multiphoton- (MP) to tunneling-ionization plus recapture (FTI)
\vspace{10pt}
\item In this paper, the followings are done
\begin{itemize}
\footnotesize
\item[$\bullet$] Experimental investigation of excitation rates of $\rm Ar$
as a function of CEP in the tunneling regime
\item[$\bullet$] Showing that bound-state populations dependence on CEP 
is strong in the tunneling regime
\item[$\bullet$] Showing at which instant between MP and tunneling regimes,
dynamics changes
\end{itemize}
\end{itemize}

\end{textblock*}

\end{frame}


\begin{frame}
\framesubtitle{Experimental setup}
\frametitle{Carrier-Envelope Phase-Dependent Strong-Field Excitation \\
\small (Phys. Rev. Lett. \textbf{128}, 173201 (2022))}

\small

\begin{textblock*}{100pt}(280pt,100pt)
\includegraphics[scale=.30]{images/medium.png}
\end{textblock*}

\begin{textblock*}{250pt}(30pt,60pt)
\begin{itemize}
\footnotesize
%\item Linear-photon-momentum transfer $\left\langle \Delta p_x\right\rangle$
%is around $4\times 10^{-4}\,\rm a.u$ for Ti:Sapphire  $(800\,\rm nm)$.
\item CEP-sytabilized laser pulses from Femto Power Pro CE-Phase 
\begin{itemize}
\footnotesize
\item[$\bullet$] $I_1 = 2.75\times 10^{14} {\rm W/cm^2}$
and $I_2 = 2.95\times 10^{14} {\rm W/cm^2}$
\item[$\bullet$] $6\,\rm fs$ pulse duration
\item [$\bullet$] (FWHM) centered at $800\,\rm nm$, 
\end{itemize}
\item Commercial $f - 2f$ interferometer (Menlo APS800)
\begin{itemize}
\footnotesize
\item[$\bullet$] Located close to the interaction region
\item[$\bullet$] Point is compensating long term drift and 
measuring the stability of the locking system
\end{itemize}

\item CEP is locked to an absolute (but not absolutely known) value 
\item Relative shift is controlled by wedge $(1\,\rm mm \sim 1.23\,\rm rad)$ 
\item Yields of $\rm Ar^*$ surviving a $0.15-0.6\,\rm ms$ flight time to 
a microchannel plate detector (MCP) measured
\item Time-of-flight used to distinguish yields of $\rm Ar^+$ 
from neutral atoms
\end{itemize}
\end{textblock*}


\begin{textblock*}{200pt}(300pt,70pt)
\tiny
Chetty \textit{et al.}, Phys. Rev. A 101, 053402 (2020)
\end{textblock*}
\end{frame}
%
%
\begin{frame}
\frametitle{Carrier-Envelope Phase-Dependent Strong-Field Excitation \\
\small (Phys. Rev. Lett. \textbf{128}, 173201 (2022))}
\framesubtitle{Experimentally measured and volume averaged TDSE results for 
the 2 sets}

\begin{textblock*}{100pt}(300pt,70pt)
\includegraphics[scale=.50, trim={0 0 250 0}, clip]{images/medium1.png} \\
\includegraphics[scale=.50, trim={250 0 0 0}, clip]{images/medium1.png} 
\end{textblock*}

\small
\begin{textblock*}{250pt}(30pt,60pt)
\begin{itemize}
\footnotesize
\item At intensity $I \leqslant 2.5\times 10^{14}\,\rm W/cm^2$, 
domination of long-term instability of the laser, thus no CEP effect observed
\item Ratio of $\rm Ar^*/Ar^+$ is observed, 
as experimental systematics are removed
\item Similar level of modulation on the order of $5\%$ and a period of 
$\pi$ is observed, as required by the symmetry
\item Experimentally measured yields are affected by 
intensity-volume averaging and CEP fluctuations
\item Numerical results have been processed to account for those fluctuations
\item Agreement between TDSE predictions and experiment shows
that numerical results are a good representation of the interaction process
\end{itemize}
\end{textblock*}



\end{frame}



 
 
\begin{frame}
\frametitle{Carrier-Envelope Phase-Dependent Strong-Field Excitation \\
\small (Phys. Rev. Lett. \textbf{128}, 173201 (2022))}
\framesubtitle{Total excitation probabilities for TDSE calculations as a 
function of intensity}

\begin{textblock*}{100pt}(300pt,70pt)
\includegraphics[scale=.30]{images/medium2.png}
\end{textblock*}

\small
\begin{textblock*}{270pt}(30pt,60pt)
\begin{itemize}
\footnotesize
\item TDSE is solved by splitting the wave function into
an inner and outer regions
\item Outer wave function is projected into the $\boldsymbol{p}-$space, 
and total ioniz. probability is obtained by integration through the space
\item Inner-region wave function is projected on the atomic states
and summed to get the probability $P({\rm Ar^*})$
\item Fig. 2(a) Comparison of $P({\rm Ar^*})$ for a cosine and sine pulses.
\begin{itemize}
\footnotesize
%\item[$\bullet$] At low intensity, $P({\rm Ar^*})$maximizes before 
%for the sine pulse
\item[$\bullet$] This behaviour changes around $I\sim 1.25\,\rm W/cm^2$,
 $\gamma \sim 1$
\item[$\bullet$] This is further edvidenced by the phase sihift 
$\Delta\sigma$ on the figure 
\item[$\bullet$] This indicates the transition
from multiphoton- to tunneling-regime
\end{itemize}
\item Fig. 2(b) $P({\rm Ar^*})$ for each intensity scaled to the the 
maximum $P({\rm Ar^*})$ for that intensity
\begin{itemize}
\footnotesize
\item[$\bullet$] Similar patterns repeat at same time as channel-closing
\item[$\bullet$] Below the 15-photons cc $(\gamma \sim 1)$, 
change is around 20\%
\item[$\bullet$] In the tunneling-regime, change is around 40\%
\end{itemize}
\end{itemize}
\end{textblock*}


\end{frame}


\begin{frame}
\framesubtitle{Details on the numerical calculations-1}
\frametitle{Carrier-Envelope Phase-Dependent Strong-Field Excitation \\
\small (Phys. Rev. Lett. \textbf{128}, 173201 (2022))}

\footnotesize
\begin{textblock*}{100pt}(240pt,100pt)
\includegraphics[scale=.30]{images/field.png}
\end{textblock*}


\begin{textblock*}{250pt}(30pt,60pt)
\begin{itemize}
\footnotesize
\item Definition of sine and cosine pulses
\begin{equation}
A(t) = \frac{E_0}{\omega}\sin^2\left(\frac{\pi t}{T}\right)
\sin\left( \omega\left(t - \frac{T}{2} + \delta_{\rm CEP}\right)\right), 
\;\;T=\frac{2\pi N}{\omega}
\nonumber
\end{equation}
\item Comparison of theoretical and experimental yields
\begin{itemize}
\footnotesize
\item[$\bullet$] Focal-volume averaging
\begin{align}
& I\left( r, z \right) = I_0
\exp\left(-\frac{2r^2}{\omega_0^2}\right), \;\;\; z_a \ll z_R 
\nonumber \\
& r\left(I^\prime\right) 
 = \omega_0 \sqrt{\frac{1}{2}\ln\left(\frac{I_0}{I^\prime}\right)}, \;\;\;
V\left( I^\prime\right) = \pi z_a r^2\left( I^\prime\right), \hspace{100pt}
\nonumber, \\
& Y_{e,i}\left( \phi, I_{peak} \right)
=\int_0^{I_{peak}} P_{e,i}\left(\phi,I\right)\frac{dV}{dI}dI, \nonumber
\end{align}
\end{itemize}
\end{itemize}
\end{textblock*}
\end{frame}




\begin{frame}
\framesubtitle{Details on the numerical calculations-2}
\frametitle{Carrier-Envelope Phase-Dependent Strong-Field Excitation \\
\small (Phys. Rev. Lett. \textbf{128}, 173201 (2022))}

\footnotesize

\begin{textblock*}{200pt}(300pt,70pt)
\tiny
Chetty \textit{et al.}, Phys. Rev. A 101, 053402 (2020) \\
Tong and Lin, J. Phys. B: At. Mol. Opt. Phys. 38 (2005)
\end{textblock*}

\begin{textblock*}{250pt}(30pt,60pt)
\begin{itemize}
\footnotesize
\item This paper is about numerical calculations that have been proved 
to be consistent with experiment
\item By TDSE calculations, they certainly mean ADK theory and not 
"numerical resolution of the TDSE"
\item Even though this is not explicitly stated in the paper, 
this is made likely by the cited references
\begin{align}
& W_{\rm TI}(F) = \frac{C_l^2}{2^{|m|}|m|!}
\frac{(2l+1)(l+|m|)!}{2(l-|m|)!}
\frac{1}{\kappa^{2Z_c/\kappa-1}}
\left( \frac{2\kappa^3}{F}\right)^{2Z_c/\kappa - |m| - 1}
e^{-2\kappa^3/3F},
\hspace{10pt} \kappa=\sqrt{2I_p} \nonumber
\end{align}
\vspace{30pt}
\begin{align}
& P = 1.0 - e^{-\int W_{\rm TI}[F(t)]dt} \nonumber
\end{align}
\end{itemize}
\end{textblock*}
\end{frame}



\begin{frame}
\frametitle{Carrier-Envelope Phase-Dependent Strong-Field Excitation \\
\small (Phys. Rev. Lett. \textbf{128}, 173201 (2022))}
\framesubtitle{Normalized probabilities for excitation to different 
$l$ states}

\begin{textblock*}{100pt}(280pt,70pt)
\includegraphics[scale=.30]{images/medium3.png}
\end{textblock*}

\small
\begin{textblock*}{250pt}(30pt,60pt)
\begin{itemize}
\footnotesize
\item In the MP picture, excitation follows the 
$\rm Ar$ dipole selection rule
\item Authors showed elsewhere that parity effect is observed for CEP-averaged
pulses
\item It is shown here that this is true for sine and cosine pulses
\item Fig. 3(a) and (b) Distribution of $l$ states due to a cosine- 
and sine- pulses
\item Fig. 3(c) Normalized probability for same parity excitated states 
\begin{itemize}
\footnotesize
\item[$\bullet$] At low intensity, parity effect is observed for both CEPs, 
consistent with expectation for MP model
\item[$\bullet$] For cosine pulses, parity is observed at all intensities
\item[$\bullet$] For sine pulses, parity conservation is lost for $\gamma > 1$
\end{itemize}
\item Venzke \textit{et al.} showed that parity 
can break down with few cycles pulses
\item \textbf{Parity depends on CEP for higher intensities}
\item Since cosine pulse doesn't break parity, 
spatial symmetry of the field matters in the tunneling-regime 
\end{itemize}
\end{textblock*}

\end{frame}




\begin{frame}
\frametitle{Carrier-Envelope Phase-Dependent Strong-Field Excitation \\
\small (Phys. Rev. Lett. \textbf{128}, 173201 (2022))}
\framesubtitle{Normalized excitation and ionization spectra}

\begin{textblock*}{100pt}(280pt,70pt)
\includegraphics[scale=.30]{images/medium4.png}
\end{textblock*}

\small
\begin{textblock*}{250pt}(30pt,60pt)
\begin{itemize}
\footnotesize
\item Quantum Defect Theory says that photo-excitation and -ionization
can be treated the same in the perturbative limit
\item Gao and Tong showed that the same holds in the nonperturbative regime
\item Fig. 4 shows the normalized photoelectrton spectra accross the ionization threshold for $80$(a) and $295\times 10^{14}\,\rm W/cm^2$(b)
\item CEP pertaining to min and max excitation yields are shown
\item (a) Lower intensity is dominated by the 13-photon resonnance ($5g$ state)
\begin{itemize}
\footnotesize
\item[$\bullet$] With CEP from $0$ to $0.6\pi$, 
$15\%$ change in modulation 
\item[$\bullet$] Change is less significant for bound states of lower energies
\end{itemize}
\item (b) For higher intensity, drastic change can be seen all along
\begin{itemize}
\footnotesize
\item[$\bullet$] Both structure and positions of the peaks are modified
with the CEP
\item[$\bullet$] Populations of individual states is extremely CEP-sensitive
\end{itemize}
\end{itemize}
\end{textblock*}

\end{frame}


